\section{Measuring Party Positions from Party Manifestos}\label{sec:manifesto-deepdive}\label{sec:manifesto-llm}

Party manifestos are the framework's sharpest available test, because
the Manifesto Project supplies the same measurement problem in two
linked data layers. The Main Dataset has one row per electoral program,
with document-level category shares, total coded quasi-sentence counts,
and derived party-position measures such as RILE; the Manifesto Corpus
stores the text layer, and for digitally annotated platforms gives the
sequence of quasi-sentences and their CMP category codes under the same
party--election identifiers
\citep{ManifestoProject2025a,ManifestoCorpus2025}. At the document
level, expert-survey means place each platform on seven-point policy
scales \citep{BenoitEtAl2025}: a trusted target that exists only at the
root of any tree. At the quasi-sentence level, trained coders assign
every clause-sized statement one category code from the Comparative
Manifesto Project scheme. Those human labels are not labels of our tree
merges. They are leaf-level statement labels; because the code counts
aggregate additively over contiguous spans, they let us induce a gold
oracle value for every leaf, every internal span, and the root. The two
levels probe different halves of the framework. With document labels
only, the local laws can enter only as audit probes and prompt-design
pressure; the question is whether a tree of local calls preserves the
expert target (Sections~\ref{sec:benoit-headline}
and~\ref{sec:rile-fg-ladder}). With induced gold labels at every node,
the local laws become training signals, and the question is whether a
summarizer trained on them actually composes
(Sections~\ref{sec:manifesto-qsentence} through~\ref{sec:gold-seeding}).
Section~\ref{sec:manifesto-lessons} collects what the corpus teaches
the framework.

\subsection{Rubrics, Surrogates, and the Measurement Problem}\label{sec:social-measurement}

Manifesto scoring is a measurement problem, not a long-context
summarization benchmark: a researcher has a rubric, a trusted but
expensive target judgment, and a downstream analysis that can be biased
if cheaper labels preserve surface accuracy while dropping
task-relevant distinctions. That is the setting in which the oracle
language of Sections~\ref{sec:framework}--\ref{sec:theorems} matters
most.

\paragraph{Rubrics, expert means, and agreement.}
Comparative-politics text measurement is built around expert-defined
constructs: economic left--right placement, social liberalism, immigration,
European integration, environmental policy, decentralization, and related
issue scales. The Manifesto Project provides the long-document corpus
\citep{ManifestoProject2025a}; \citet{BenoitEtAl2025} provide the recent
LLM benchmark we use, comparing model-produced scores to expert-survey
means on seven-point policy scales. In this design, the target is a
construct-level party-position estimate obtained by aggregating expert
judgments, not a single annotator's sentence label; split-expert
correlations are an empirical agreement reference for how noisy that
measurement target is.

This matters for C-TreePO because the object being preserved is the
rubric-indexed judgment, not the literal text. A summary can safely discard
stylistic details, slogans, and redundant examples if the expert-scale
position remains unchanged. It fails when it drops evidence that changes
the issue placement or a cross-section interaction that an expert would
have used to resolve an ambiguous platform.

\paragraph{Surrogate labels and downstream inference.}
LLMs are attractive in this setting because they can cheaply produce
rubric labels, explanations, and summaries for many long documents.
Social-science use demands more than correlation with expert means.
\citet{EgamiEtAl2024} show that using imperfect surrogate
labels in downstream inference can induce bias even when prediction
quality is high, and their design-based supervised-learning framework
uses gold-standard annotations to correct that risk. C-TreePO addresses
the complementary compression question: when a long document is replaced
by a tree summary, what evidence has to survive for the downstream
oracle-indexed claim to remain valid?

The audit in Section~\ref{sec:estimation} is therefore a measurement
instrument, not a generic quality score. It samples leaf, idempotence, and
merge checks from the realized tree; the calibration residual measures the
gap between the judge used during training or auditing and the expert
oracle the researcher ultimately trusts; and the finite-sample envelope
states how much uncertainty remains after collecting global and local
labels.

\paragraph{Why the Manifesto benchmark fits.}
Manifestos have exactly the structure the framework is meant to test.
They are long, sectioned documents; their issue positions are rubric-coded;
their evidence is often local but can be resolved by interactions across
sections; and the benchmark reports external expert means rather than only
model-internal consistency. The linked Main Dataset/Corpus design adds a
rare diagnostic feature: for annotated platforms, we observe both the
full-document MPDS rollup and the quasi-sentence code sequence that
produces it. We can therefore ask not only whether a full tree predicts
the document label, but also where local compression helps or fails along
the way. The interpretation to carry into the numbers below: matching
the benchmark is evidence that a single open-weight, tree-structured
prompt program can preserve a social-science measurement target, while
the local-law audit supplies the failure modes and finite-sample caveats
that a flat long-context read does not expose.
